\section{Peluang Usaha dalam Berwirausaha}
\subsection{Membangun Usaha (Berwirausaha)}
\begin{enumerate}
\item Wirausaha merupakan konsep beraktivitas mandiri dengan mengandalkan kekuatan diri sendiri untuk berinovasi dan berkreasi untuk menghasilkan nilai ekonomi dalam pemenuhan kebutuhan hidup pelanggan, pelaku dalam wirausaha disebut wirausahawan, yang memiliki karakteristik tertentu yang memengaruhi keberhasilan usaha.
\item Motivasi Berwirausaha:
  \begin{enumerate}
  \item Kemandirian dan kebebasan.
  \item Potensi keuntungan finansial.
  \item Fleksibilitas waktu.
  \item Inovasi dan kreativitas.
  \item Perubahan karier.
  \item Tantangan.
  \end{enumerate}
\item Perbedaan Pekerja dan Wirausaha:
  \begin{enumerate}
  \item Tanggung Jawab dan Kontrol:
    \begin{enumerate}
    \item Wirausaha: Memiliki dan mengendalikan bisnisnya.
    \item Pekerja: Tanggung jawabnya terbatas pada tugas dan fungsi yang ditentukan oleh perannya dalam organisasi.
    \end{enumerate}
  \item Risiko:
    \begin{enumerate}
    \item Wirausaha: Mengambil risiko yang jauh lebih besar.
    \item Pekerja: Risiko terbatas pada kehilangan pekerjaan.
    \end{enumerate}
  \item Kebebasan dan Fleksibilitas:
    \begin{enumerate}
    \item Wirausaha: Memiliki kebebasan dalam mengatur jadwal kerja dan keputusan bisnis.
    \item Pekerja: Jadwal kerja dan tugasnya lebih terstruktur dan ditentukan oleh kebijakan organisasi.
    \end{enumerate}
  \item Pendapatan:
    \begin{enumerate}
    \item Wirausaha: Pendapatan tidak tetap dan bergantung pada keberhasilan bisnis.
    \item Pekerja: Biasanya menerima pendapatan tetap berupa gaji atau upah yang disepakati, terlepas dari kinerja perusahaan.
    \end{enumerate}
  \item Pengembangan Karier:
    \begin{enumerate}
    \item Wirausaha: Karier berkembang berdasarkan pertumbuhan dan keberhasilan bisnis yang dibangunnya.
    \item Pekerja: Pengembangan karier bergantung pada struktur dan peluang dalam organisasi tempat bekerja.
    \end{enumerate}
  \item Pengambilan Keputusan:
    \begin{enumerate}
    \item Wirausaha: Berperan langsung dalam menciptakan lapangan kerja, inovasi, dan menggerakan ekonomi.
    \item Pekerja: Berkontribusi pada ekonomi melalui keterampilan dan tenaga kerja pekerja dalam berbagai industri dan sektor.
    \end{enumerate}
  \end{enumerate}
\end{enumerate}

\subsection{Potensi Keberhasilan Berwirausaha di Bidang TI}
\begin{enumerate}
\item Permintaan tinggi dan pertumbuhan cepat.
\item Potensi keuntungan yang besar.
\item Fleksibilitas dalam bekerja.
\item Kepuasan dalam inovasi dan penciptaan.
\item Dampak sosial dan perubahan.
\item Kesempatan untuk belajar dan berkembang.
\item Rendahnya hambatan masuk.
\item Komunitas dan dukungan.
\item Personalisasi produk dan layanan.
\end{enumerate}
